\chapter*{Abstract}

% PLACEHOLDER: Manter sincronizado com Cap00/Resumo.tex.
% Os trechos de RESULTS e CONCLUSION acompanham os placeholders do resumo.

Stress and anxiety are frequent complaints among young adults and university students, and brief, accessible, and enjoyable interventions for their short-term attenuation remain an open problem. This work presents the development of \textit{Cozy Pong}, a virtual reality (VR) game that combines table-tennis and rhythm-game mechanics, set under the influence of \textit{lo-fi} music, and proposes a protocol to empirically validate its effect on the regulation of stress and anxiety. The game was developed in the Unity engine for the Meta Quest 3S device and organizes gameplay around rhythmic maps (\textit{beatmaps}) synchronized to the music, with visual, auditory, and haptic feedback. Four conditions operationalize the experimental manipulation: \textit{Cozy Pong} in a relaxing configuration, the intervention under evaluation; the same game in a lively configuration, which isolates the effect of the deliberately relaxing characterization; listening to the same \textit{lo-fi} soundtrack while seated, without playing, which holds the musical stimulus constant and isolates the effect of the game; and a guided meditation session on the same VR device, a comparator representing today's established digital relaxation practice. The protocol adopts a randomized within-subject design counterbalanced by a Williams square, with thirty-two participants, each condition preceded by a standardized mental-arithmetic stress induction. The effect is measured by cross-referencing psychophysiological measures, heart-rate variability (HRV) from an electrode chest strap and electrodermal activity (EDA) from a raw-signal sensor, with a psychometric battery composed entirely of freely available scales validated in Brazilian Portuguese. The central hypothesis is that practicing light, rhythmic exercise synchronized with relaxing music promotes a more significant reduction in stress and anxiety than passively listening to the same music and than a more demanding configuration of the game itself, while matching guided meditation. \ph{summarize the main results: variations in the psychological scales and physiological indices per condition.} \ph{state the main conclusion after analysis.} The study is expected to contribute empirical evidence, at the intersection of computing and health, on the effect of playful VR interventions on short-term emotional recovery.

\vspace{1em}
\noindent\textbf{\textit{Keywords}:~} virtual reality. rhythm game. exergame. stress reduction. anxiety. lo-fi music. heart-rate variability. electrodermal activity.
